\section{Conclusion}

Even though it is well known that the OBD depends on the choice of reference group, we show that this dependence can be more consequential, leading to sign reversals.
We present a simple healthcare example, inspired by real-world data, in which one OBD direction suggests that quality of care improves patient outcomes, while the other suggests it worsens them. 
We further prove that such sign reversals are not isolated phenomena, but can occur in a wide variety of scenarios. 
Practitioners should therefore be cautious when using the OBD, at least in settings where no reference group is clearly justified. 